%% Section S9, the last section of the supplement: run provenance, hardware, configurations and reproducibility.
%% Numbers from the runs are \Res keys (results-values.tex, generated by tools/build_paper_results.py). Values that no
%% metric file holds (training configuration files, architecture constants) are typed from neural-*/outputs/code and
%% the notebook sources. The section is not wrapped in the appendices environment: sn-jnl would letter this part
%% (Appendix A, Table A1), which breaks the S numbering used by the rest of the supplement.

\section{Runs, hardware and reproducibility}\label{sec:app}

This section records what produced each saved result, on which hardware and with which time budget, and how to regenerate every number and both documents. Three notebooks produced the outputs. The data pipeline (\texttt{neural-data-pipeline}) builds the leakage-safe dataset manifest. The temporal super-resolution notebook (TSR, \texttt{neural-supersampling}) and the frame-generation notebook (FG, \texttt{neural-frame-generation}) read that manifest, train their networks and evaluate them. The saved outputs of each notebook come from one run, the run listed in \cref{tab:app-runs}.

\subsection{Run provenance}\label{sec:app-prov}

Each notebook keeps its own clock. The clock starts when the notebook creates its run object in the setup cell, and every timed stage prints the elapsed time on its line of the log. \Cref{tab:app-runs} gives the dates of each run and two durations. The summary value is written to \texttt{summary.json} at the end of the notebook. The log value is the total printed on the last timed stage, which the generator stores as \texttt{run.<notebook>.hours}.

\begin{table}[!htbp]
\caption{Dates and wall-clock hours of the three runs. The difference is the summary value minus the log value.}\label{tab:app-runs}
\footnotesize
\begin{tabular*}{\textwidth}{@{\extracolsep\fill}lllrrr}
\toprule
Notebook & Start & End & Summary (h) & Log (h) & Difference (h) \\
\midrule
Data pipeline & \Res{run.neural-data-pipeline.start} & \Res{run.neural-data-pipeline.end} & \Res{data.summary.run_hours} & \Res{run.neural-data-pipeline.hours} & \ResDiff{data.summary.run_hours}{run.neural-data-pipeline.hours} \\
TSR & \Res{run.neural-supersampling.start} & \Res{run.neural-supersampling.end} & \Res{tsr.summary.run_hours} & \Res{run.neural-supersampling.hours} & \ResDiff{tsr.summary.run_hours}{run.neural-supersampling.hours} \\
FG & \Res{run.neural-frame-generation.start} & \Res{run.neural-frame-generation.end} & \Res{fg.summary.run_hours} & \Res{run.neural-frame-generation.hours} & \ResDiff{fg.summary.run_hours}{run.neural-frame-generation.hours} \\
\botrule
\end{tabular*}
\end{table}

The two durations differ because of the code. The summary cell reads the clock after the last timed stage, and the cells between that stage and the summary are untimed. In TSR these cells aggregate the results, compute the perceptual and static-scene metrics, draw the visual comparisons, tabulate the ablations, time the latency of the fused tiers and export the ONNX models. In FG they aggregate the results, compute the UI, temporal-consistency, visual, ablation, latency and pacing results, and export the ONNX models. The untimed time accounts for the differences in \cref{tab:app-runs}. The data pipeline writes its summary straight after its last timed stage, so its two values agree. For all three runs the log value agrees to within 0.01 h with the span from the first to the last log line, and the generator takes the run dates from the same log.

The configuration cells set a time budget of \Res{tsr.config.budget_h} h for TSR and \Res{fg.config.budget_h} h for FG (\texttt{RUN\_BUDGET\_H}). The notebooks describe the budget as Kaggle's 12-hour session limit minus a 72-minute margin. That limit is the notebooks' own assumption, which this supplement does not check. Both runs finished below their budgets on either measure. The data pipeline has no budget and finished in under one hour. Neither duration includes the cells that delete the training caches and package the outputs, which run after the summary.

The executed notebooks show which code ran. The data pipeline and TSR copies have the same cell sources as a fresh build of the current \texttt{notebooks} and \texttt{src/neuralss} files, so the library code of those runs is the current code. The FG copy differs in one plotting cell, the figure of the real-suite results. That cell was changed after the run, so that it hides a panel with no data and rotates the axis labels by 20 degrees. The change affects only \texttt{real\_results.png}, and no metric file depends on it.

\subsection{Hardware and software}\label{sec:app-hw}

\Cref{tab:app-hw} gives the hardware that each notebook detected at run time. Each run found two identical Tesla T4 GPUs. The training notebooks use both GPUs together for the main stages and one GPU per process for the parallel pairs and the ablations.

\begin{table}[!htbp]
\caption{Hardware and software recorded by each run (\texttt{metrics/hardware.json})}\label{tab:app-hw}
\begin{tabular*}{\textwidth}{@{\extracolsep\fill}lccc}
\toprule
Item & Data pipeline & TSR & FG \\
\midrule
GPUs & \Res{data.hardware.gpus} \Res{data.hardware.gpu_list.0.name} & \Res{tsr.hardware.gpus} \Res{tsr.hardware.gpu_list.0.name} & \Res{fg.hardware.gpus} \Res{fg.hardware.gpu_list.0.name} \\
Memory per GPU (GB) & \Res{data.hardware.gpu_list.0.vram_gb} & \Res{tsr.hardware.gpu_list.0.vram_gb} & \Res{fg.hardware.gpu_list.0.vram_gb} \\
Compute capability & \Res{data.hardware.gpu_list.0.cc} & \Res{tsr.hardware.gpu_list.0.cc} & \Res{fg.hardware.gpu_list.0.cc} \\
Multiprocessors per GPU & \Res{data.hardware.gpu_list.0.sm} & \Res{tsr.hardware.gpu_list.0.sm} & \Res{fg.hardware.gpu_list.0.sm} \\
CPU cores & \Res{data.hardware.cpu_count} & \Res{tsr.hardware.cpu_count} & \Res{fg.hardware.cpu_count} \\
RAM (GB) & \Res{data.hardware.ram_total_gb} & \Res{tsr.hardware.ram_total_gb} & \Res{fg.hardware.ram_total_gb} \\
Free disk in the working folder (GB) & \Res{data.hardware.disk_free_gb_working} & \Res{tsr.hardware.disk_free_gb_working} & \Res{fg.hardware.disk_free_gb_working} \\
Python & \Res{data.hardware.python} & \Res{tsr.hardware.python} & \Res{fg.hardware.python} \\
PyTorch & \Res{data.hardware.torch} & \Res{tsr.hardware.torch} & \Res{fg.hardware.torch} \\
CUDA & \Res{data.hardware.cuda} & \Res{tsr.hardware.cuda} & \Res{fg.hardware.cuda} \\
\botrule
\end{tabular*}
\end{table}

The data pipeline ran on an older software image than the two training notebooks. It used PyTorch \Res{data.hardware.torch} with CUDA \Res{data.hardware.cuda}, while TSR and FG used PyTorch \Res{tsr.hardware.torch} with CUDA \Res{tsr.hardware.cuda}. The training notebook states that the library supports PyTorch 2.0 and later.

The training notebooks install, when they are missing, three optional packages at run time: \texttt{lpips}, \texttt{flip-evaluator} and \texttt{onnxruntime}. Their logs record only whether each package could be imported afterwards. All three succeeded in both training runs, but the outputs do not record package versions. The versions of the other libraries, such as NumPy and pandas, are not recorded either.

\subsection{Stage times}\label{sec:app-stage}

\FloatBarrier

Each notebook times its stages with a context manager, which records the duration of each stage in \texttt{metrics/stage\_times.csv} and draws the bar chart in \texttt{plots/stage\_times.png}. \Cref{fig:app-stage} shows the three charts, panel (\ref{sub@fig:app-stage-data}) for the data pipeline, panel (\ref{sub@fig:app-stage-tsr}) for TSR and panel (\ref{sub@fig:app-stage-fg}) for FG, and \cref{tab:app-stages} lists the values. The setup, configuration and packaging cells are not timed.

\begin{figure}[!htbp]
\caption{Minutes per timed stage, as drawn by each notebook. Panel (a) uses a logarithmic axis, panels (b) and (c) a linear one.}\label{fig:app-stage}
\begin{subfigure}{\textwidth}
\centering
\includegraphics[width=0.8\textwidth,height=0.15\textheight,keepaspectratio]{\DataOut/plots/stage_times.png}
\caption{Data pipeline}\label{fig:app-stage-data}
\end{subfigure}
\par\medskip
\begin{subfigure}{\textwidth}
\centering
\includegraphics[width=0.8\textwidth,height=0.15\textheight,keepaspectratio]{\TsrOut/plots/stage_times.png}
\caption{TSR}\label{fig:app-stage-tsr}
\end{subfigure}
\par\medskip
\begin{subfigure}{\textwidth}
\centering
\includegraphics[width=0.8\textwidth,height=0.15\textheight,keepaspectratio]{\FgOut/plots/stage_times.png}
\caption{FG}\label{fig:app-stage-fg}
\end{subfigure}
\end{figure}

\FloatBarrier

{\setlength{\LTcapwidth}{\textwidth}
\begin{longtable}{lrrr}
\caption{Minutes per timed stage of each run (\texttt{metrics/stage\_times.csv}). A dash means that the notebook has no such stage.}\label{tab:app-stages}\\
\toprule
Stage & Data pipeline & TSR & FG \\
\midrule
kaggle\_discovery & \Res{data.stage_times.kaggle_discovery.minutes} & -- & -- \\
hf\_materialisation & \Res{data.stage_times.hf_materialisation.minutes} & -- & -- \\
keyframe\_stats & \Res{data.stage_times.keyframe_stats.minutes} & -- & -- \\
filters & \Res{data.stage_times.filters.minutes} & -- & -- \\
hashes & \Res{data.stage_times.hashes.minutes} & -- & -- \\
embeddings & \Res{data.stage_times.embeddings.minutes} & -- & -- \\
local\_features & \Res{data.stage_times.local_features.minutes} & -- & -- \\
calibration & \Res{data.stage_times.calibration.minutes} & -- & -- \\
duplicate\_search & \Res{data.stage_times.duplicate_search.minutes} & -- & -- \\
split\_and\_audit & \Res{data.stage_times.split_and_audit.minutes} & -- & -- \\
manifest & \Res{data.stage_times.manifest.minutes} & \Res{tsr.stage_times.manifest.minutes} & \Res{fg.stage_times.manifest.minutes} \\
caches & -- & \Res{tsr.stage_times.caches.minutes} & \Res{fg.stage_times.caches.minutes} \\
batch\_probe & -- & \Res{tsr.stage_times.batch_probe.minutes} & \Res{fg.stage_times.batch_probe.minutes} \\
train\_teacher & -- & \Res{tsr.stage_times.train_teacher.minutes} & \Res{fg.stage_times.train_teacher.minutes} \\
train\_student & -- & \Res{tsr.stage_times.train_student.minutes} & \Res{fg.stage_times.train_student.minutes} \\
train\_pair & -- & \Res{tsr.stage_times.train_pair.minutes} & \Res{fg.stage_times.train_pair.minutes} \\
ablations & -- & \Res{tsr.stage_times.ablations.minutes} & \Res{fg.stage_times.ablations.minutes} \\
test\_suites & -- & \Res{tsr.stage_times.test_suites.minutes} & \Res{fg.stage_times.test_suites.minutes} \\
evaluation & -- & \Res{tsr.stage_times.evaluation.minutes} & \Res{fg.stage_times.evaluation.minutes} \\
\botrule
\end{longtable}
}

\FloatBarrier

In \cref{fig:app-stage-data} the Hugging Face transfer and the keyframe decoding are the two longest bars, and the manifest stage is too short to see. In \cref{fig:app-stage-tsr} and \cref{fig:app-stage-fg} the four training stages are the four longest bars. The student stage is the longest bar in both training runs, and the caches stage is the longest stage that is not a training stage.

Training dominates the time budget of both training runs. The four training stages take longer than all the other stages together, and the student stage is the single largest stage. The caches and the evaluation are short by comparison, and the manifest is negligible in both. \Cref{tab:app-budget} compares each budgeted stage with its plan. The FG run had a caches stage longer than its planned cap. The notebook then shrank every training budget slightly, before training started, so that the run stayed within its ceiling. The data pipeline has no training, and its Hugging Face transfer and keyframe decoding together take longer than its other nine stages combined.

\begin{table}[!htbp]
\caption{Planned and measured minutes of the budgeted stages. The plan is the configuration value. The FG plan is scaled slightly down before training, so its measured times are close to the scaled plan, not equal to the configuration value.}\label{tab:app-budget}
\begin{tabular*}{\textwidth}{@{\extracolsep\fill}lrrrr}
\toprule
Stage & TSR plan & TSR measured & FG plan & FG measured \\
\midrule
Caches (at most) & \Res{tsr.config.cfg.cache.build_min} & \Res{tsr.stage_times.caches.minutes} & \Res{fg.config.cfg.cache.build_min} & \Res{fg.stage_times.caches.minutes} \\
Teacher & \Res{tsr.config.cfg.minutes.teacher} & \Res{tsr.stage_times.train_teacher.minutes} & \Res{fg.config.cfg.minutes.teacher} & \Res{fg.stage_times.train_teacher.minutes} \\
Student & \Res{tsr.config.cfg.minutes.student} & \Res{tsr.stage_times.train_student.minutes} & \Res{fg.config.cfg.minutes.student} & \Res{fg.stage_times.train_student.minutes} \\
Pair of parallel runs & \Res{tsr.config.cfg.minutes.pair} & \Res{tsr.stage_times.train_pair.minutes} & \Res{fg.config.cfg.minutes.pair} & \Res{fg.stage_times.train_pair.minutes} \\
Ablations (three pairs) & $3\times$\Res{tsr.config.cfg.minutes.ablation} & \Res{tsr.stage_times.ablations.minutes} & $3\times$\Res{fg.config.cfg.minutes.ablation} & \Res{fg.stage_times.ablations.minutes} \\
\botrule
\end{tabular*}
\end{table}

\FloatBarrier

\subsection{Configuration}\label{sec:app-cfg}

Each notebook writes its configuration to \texttt{metrics/config.json}. The file holds the configuration cell's dictionary, the run budget and, for the training notebooks, the ablation overrides. The tables below cover every key that affects a run. Their values are generated keys, so a new run updates them. Keys that exist only in the training configuration files (\texttt{code/cfg\_*.json}) are typed, because no metric file holds them.

\subsubsection{Data pipeline}\label{sec:app-cfg-data}

\Cref{tab:app-cfg-data} lists the keys of the data pipeline configuration, which sets the caps of each source, the copy-detection thresholds and the split.

\begin{longtable}{@{}p{0.31\textwidth}p{0.12\textwidth}p{0.45\textwidth}@{}}
\caption{Data pipeline configuration. Keys are under \texttt{cfg} in \texttt{metrics/config.json}.}\label{tab:app-cfg-data}\\
\toprule
Key & Value & Meaning \\
\midrule
\endfirsthead
\toprule
Key & Value & Meaning \\
\midrule
\endhead
\midrule
\multicolumn{3}{r}{\textit{Continued on the next page}} \\
\endfoot
\botrule
\endlastfoot
\multicolumn{3}{@{}l}{\textit{Discovery and seeds}} \\
\texttt{seed} & \Res{data.config.cfg.seed} & seed of the randomised discovery walk and of the split hash \\
\texttt{list\_cap\_files} & \Res{data.config.cfg.list_cap_files} & most files listed per Kaggle dataset (a source may set its own cap) \\
\texttt{workers} & \Res{data.config.cfg.workers} & worker processes for image decoding and SIFT features \\
\multicolumn{3}{@{}l}{\textit{Hugging Face subsets}} \\
\texttt{hf\_caps.game\_stills} & \Res{data.config.cfg.hf_caps.game_stills.0} / \Res{data.config.cfg.hf_caps.game_stills.1} & game still images for training and for test \\
\texttt{hf\_caps.tartanair} & \Res{data.config.cfg.hf_caps.tartanair} & frames per TartanAir environment \\
\texttt{hf\_caps.tartanair\_seq} & \Res{data.config.cfg.hf_caps.tartanair_seq} & frames per TartanAir sequence \\
\texttt{hf\_caps.gameir} & \Res{data.config.cfg.hf_caps.gameir.0} / \Res{data.config.cfg.hf_caps.gameir.1} & GameIR clips for training and for test \\
\texttt{hf\_caps.vimeo1080p} & \Res{data.config.cfg.hf_caps.vimeo1080p.0} / \Res{data.config.cfg.hf_caps.vimeo1080p.1} & Vimeo-1080p videos for training, and the videos of the pinned validation shard, which become test items \\
\multicolumn{3}{@{}l}{\textit{Copy detection and split}} \\
\texttt{calibration\_pairs} & \Res{data.config.cfg.calibration_pairs} & synthetic copy pairs for calibrating the duplicate thresholds \\
\texttt{dedup.k} & \Res{data.config.cfg.dedup.k} & embedding neighbours per keyframe that become candidate copies \\
\texttt{dedup.retrieval\_\allowbreak recall} & \Res{data.config.cfg.dedup.retrieval_recall} & share of synthetic copies that the embedding threshold must retrieve \\
\texttt{dedup.sift\_features} & \Res{data.config.cfg.dedup.sift_features} & SIFT keypoints per image in geometric verification \\
\texttt{dedup.min\_std\_hash} & \Res{data.config.cfg.dedup.min_std_hash} & grey-level standard deviation a keyframe needs before pHash matching \\
\texttt{dedup.max\_hash\_\allowbreak candidates\_per\_kf} & \Res{data.config.cfg.dedup.max_hash_candidates_per_kf} & the pHash threshold is tightened until at most this many candidate pairs per keyframe remain, or until the exact threshold is reached \\
\texttt{dedup.audit\_k} & \Res{data.config.cfg.dedup.audit_k} & nearest reference keyframes checked for each validation or test keyframe \\
\texttt{dedup.audit\_rounds} & \Res{data.config.cfg.dedup.audit_rounds} & most audit rounds, each adding the duplicates it finds as edges and splitting again \\
\texttt{val\_fraction} & \Res{data.config.cfg.val_fraction} & share of train and validation components assigned to validation, by a hash of the component root \\
\end{longtable}

\FloatBarrier

The sources are the datasets that the manifest draws on. \Cref{tab:app-sources} lists them as \texttt{config.json} records them. Every source was reachable in both training runs.

\begin{table}[!htbp]
\caption{Sources of the manifest. A host of \texttt{kaggle} means a Kaggle dataset and \texttt{hf} a Hugging Face repository. The kinds are single images (still), videos, rendered sequences (render\_seq), rendered image pairs (render\_pairs) and video files (video\_file).}\label{tab:app-sources}
\footnotesize
\begin{tabular*}{\textwidth}{@{\extracolsep\fill}lllp{0.42\textwidth}}
\toprule
Source & Host & Kind & Identifier \\
\midrule
DIV2K & Kaggle & still & \texttt{soumikrakshit/\allowbreak div2k-high-resolution-images} \\
Flickr2K & Kaggle & still & \texttt{daehoyang/\allowbreak flickr2k} \\
REDS & Kaggle & video & \texttt{amithkesavmrajagiri/\allowbreak reds-dataset} \\
Vimeo-90K septuplets & Kaggle & video & \texttt{wangsally/\allowbreak vimeo-90k-7} \\
Vimeo-90K triplets & Kaggle & video & \texttt{chenshu123/\allowbreak vimeo-triplet} \\
MPI Sintel & Kaggle & render\_seq & \texttt{artemmmtry/\allowbreak mpi-sintel-dataset} \\
Vid4 & Kaggle & video & \texttt{uom200647r/\allowbreak vid4-dataset} \\
Game stills & Hugging Face & still & \texttt{ericphann/\allowbreak video-game-super-resolution} \\
TartanAir & Hugging Face & render\_seq & \texttt{theairlabcmu/\allowbreak tartanair} \\
GameIR & Hugging Face & render\_pairs & \texttt{LLLebin/\allowbreak GameIR} \\
Vimeo-1080p & Hugging Face & video\_file & \texttt{danjacobellis/\allowbreak vimeo1080p} \\
\botrule
\end{tabular*}
\end{table}

\subsubsection{Temporal super-resolution}\label{sec:app-cfg-tsr}

The TSR training runs use AdamW \cite{loshchilov2019adamw} and FP16 autocast with loss scaling \cite{micikevicius2018mixed}. \Cref{tab:app-cfg-tsr} lists the keys of the TSR configuration. The last group holds the keys that appear only in the training configuration files, and their values are typed.

The TSR tiers differ in the width and depth of the trunk and in the built-in motion estimator. \Cref{tab:app-tsr-tiers} lists them. The trunk is a stack of RepConv blocks \cite{ding2021repvgg} and the output is unfolded with pixel shuffle \cite{shi2016espcn}. The GMAC column gives billions of multiply-accumulate operations (MAC) per 1080p output frame, counted on the fused network.

{\footnotesize\setlength{\LTcapwidth}{\textwidth}
\begin{longtable}{lrrrrlrrr}
\caption{TSR tiers. Width, blocks and U-Net blocks describe the trunk, and hidden is the number of channels of the recurrent state. The motion estimator is given as its widths and its residual blocks, and the refinement is the width of the output-resolution head (0 means none). The parameters are the training-time counts, and the cost is per 1080p output frame at 3x.}\label{tab:app-tsr-tiers}\\
\toprule
Tier & Width & Blocks & U-Net & Hidden & Motion & Refine & Params & GMAC \\
\midrule
\tier{igpu} & 16 & 3 & 0 & 8 & 32/24/16, 2 & 0 & \Res{tsr.model_tiers.igpu.3.train_params} & \Res{tsr.model_tiers.igpu.3.gmac_per_frame_1080p} \\
\tier{low} & 24 & 4 & 2 & 8 & 48/32/24, 3 & 8 & \Res{tsr.model_tiers.low.3.train_params} & \Res{tsr.model_tiers.low.3.gmac_per_frame_1080p} \\
\tier{mid} & 32 & 4 & 3 & 16 & 64/48/32, 4 & 8 & \Res{tsr.model_tiers.mid.3.train_params} & \Res{tsr.model_tiers.mid.3.gmac_per_frame_1080p} \\
\tier{high} & 48 & 6 & 4 & 16 & 96/64/48, 4 & 12 & \Res{tsr.model_tiers.high.3.train_params} & \Res{tsr.model_tiers.high.3.gmac_per_frame_1080p} \\
\tier{teacher} & 64 & 8 & 6 & 32 & 128/96/64, 6 & 16 & \Res{tsr.model_tiers.teacher.3.train_params} & \Res{tsr.model_tiers.teacher.3.gmac_per_frame_1080p} \\
\botrule
\end{longtable}
}

\FloatBarrier

\begin{longtable}{@{}p{0.31\textwidth}p{0.12\textwidth}p{0.45\textwidth}@{}}
\caption{TSR configuration. Keys are under \texttt{cfg} in \texttt{metrics/config.json} unless shown otherwise.}\label{tab:app-cfg-tsr}\\
\toprule
Key & Value & Meaning \\
\midrule
\endfirsthead
\toprule
Key & Value & Meaning \\
\midrule
\endhead
\midrule
\multicolumn{3}{r}{\textit{Continued on the next page}} \\
\endfoot
\botrule
\endlastfoot
\multicolumn{3}{@{}l}{\textit{Budget, model and data}} \\
\texttt{budget\_h} (top level) & \Res{tsr.config.budget_h} & ceiling of the run in hours \\
\texttt{seed} & \Res{tsr.config.cfg.seed} & seed of the training run \\
\texttt{scale} & \Res{tsr.config.cfg.scale} & primary super-resolution factor, from 360p to 1080p \\
\texttt{scale\_alt} & \Res{tsr.config.cfg.scale_alt} & secondary factor, used by the 2x student \\
\texttt{patch} & \Res{tsr.config.cfg.patch} & output patch in pixels. At 3x, the low-resolution patch has one third of this side \\
\texttt{seq\_len} & \Res{tsr.config.cfg.seq_len} & frames per training sequence, the length of back-propagation through time \\
\texttt{still\_size} & \Res{tsr.config.cfg.still_size} & side in pixels of the texture crops from which synthetic scenes are rendered \\
\texttt{student\_tier}, \texttt{small\_tier}, \texttt{teacher\_tier} & low, igpu, teacher & tiers of the student, of the integrated-GPU run and of the teacher (\cref{tab:app-tsr-tiers}) \\
\multicolumn{3}{@{}l}{\textit{Caches (requested counts)}} \\
\texttt{still\_train}, \texttt{still\_val} & \Res{tsr.config.cfg.cache.still_train} / \Res{tsr.config.cfg.cache.still_val} & still-image crops for training and for validation \\
\texttt{win\_train}, \texttt{win\_val} & \Res{tsr.config.cfg.cache.win_train} / \Res{tsr.config.cfg.cache.win_val} & sequence windows of \texttt{seq\_len} frames for training and for validation \\
\texttt{build\_min} & \Res{tsr.config.cfg.cache.build_min} & total minutes allowed for decoding all caches \\
\multicolumn{3}{@{}l}{\textit{Sampling}} \\
\texttt{p\_real} & \Res{tsr.config.cfg.p_real} & share of each batch drawn from real sequences, the rest being synthetic \\
\texttt{p\_video\_synth}, \texttt{p\_video\_real} & \Res{tsr.config.cfg.p_video_synth} / \Res{tsr.config.cfg.p_video_real} & share of synthetic and of real samples turned into plain video (no jitter, depth or motion vectors, with 8-bit quantisation and mild degradation) \\
\texttt{p\_video\_mv} & \Res{tsr.config.cfg.p_video_mv} & share of those plain-video samples that keep estimated motion vectors \\
\texttt{p\_drop\_mv}, \texttt{p\_drop\_depth} & \Res{tsr.config.cfg.p_drop_mv} / \Res{tsr.config.cfg.p_drop_depth} & share of render-like samples whose motion vectors or depth are withheld, so that the network also learns without them \\
\texttt{long\_every} & \Res{tsr.config.cfg.long_every} & every this many steps trains on a long unroll of synthetic frames \\
\texttt{long\_T}, \texttt{long\_patch} & \Res{tsr.config.cfg.long_t} / \Res{tsr.config.cfg.long_patch} & frames and patch size of a long unroll \\
\texttt{mv\_warmup\_frac} & \Res{tsr.config.cfg.mv_warmup_frac} & share of the run during which samples without motion vectors warp the history with the label motion, with a probability that falls from one to zero \\
\multicolumn{3}{@{}l}{\textit{Optimisation}} \\
\texttt{lr}, \texttt{lr\_teacher} & \Res{tsr.config.cfg.lr} / \Res{tsr.config.cfg.lr_teacher} & peak learning rate of the students and of the teacher, followed by warm-up and a cosine decay over the wall-clock budget \\
\texttt{min\_lr} & \Res{tsr.config.cfg.min_lr} & floor of the cosine decay \\
\texttt{warmup\_frac} & \Res{tsr.config.cfg.warmup_frac} & share of the budget spent warming the learning rate up \\
\texttt{weight\_decay} & \Res{tsr.config.cfg.weight_decay} & AdamW weight decay \\
\texttt{ema\_decay} & \Res{tsr.config.cfg.ema_decay} & decay of the exponential moving average of the weights, which is the copy that is evaluated \\
\texttt{grad\_clip} & \Res{tsr.config.cfg.grad_clip} & gradient norm above which the gradients are scaled down \\
\multicolumn{3}{@{}l}{\textit{Loss weights (per-frame loss terms)}} \\
\texttt{w\_grad}, \texttt{w\_fft} & \Res{tsr.config.cfg.w_grad} / \Res{tsr.config.cfg.w_fft} & weights of the gradient loss and of the loss on Fourier magnitudes \\
\texttt{w\_temporal} & \Res{tsr.config.cfg.w_temporal} & weight of the temporal-change loss along known motion \\
\texttt{w\_flow} & \Res{tsr.config.cfg.w_flow} & weight of the motion-estimation loss, per low-resolution pixel \\
\texttt{w\_lpips} & \Res{tsr.config.cfg.w_lpips} & weight of the LPIPS perceptual loss, set to zero when the package is missing \\
\texttt{w\_distill} & \Res{tsr.config.cfg.w_distill} & weight of distillation from the teacher, zero because no teacher is used \\
\multicolumn{3}{@{}l}{\textit{Validation and budgets}} \\
\texttt{val\_bs}, \texttt{val\_batches} & \Res{tsr.config.cfg.val_bs} / \Res{tsr.config.cfg.val_batches} & samples per validation batch and batches per validation round \\
\texttt{val\_interval\_min} & \Res{tsr.config.cfg.val_interval_min} & longest time between validation rounds, in minutes \\
\texttt{minutes.teacher}, \texttt{minutes.student} & \Res{tsr.config.cfg.minutes.teacher} / \Res{tsr.config.cfg.minutes.student} & wall-clock training budget of the teacher and of the student, in minutes \\
\texttt{minutes.pair}, \texttt{minutes.ablation} & \Res{tsr.config.cfg.minutes.pair} / \Res{tsr.config.cfg.minutes.ablation} & budget of each parallel pair, and of each ablation pair, in minutes \\
\multicolumn{3}{@{}l}{\textit{Evaluation}} \\
\texttt{synthetic\_per\_class} & \Res{tsr.config.cfg.eval.synthetic_per_class} & synthetic test crops per motion class, of four classes \\
\texttt{synthetic\_T}, \texttt{synthetic\_P} & \Res{tsr.config.cfg.eval.synthetic_t} / \Res{tsr.config.cfg.eval.synthetic_p} & frames and output size of each synthetic test sequence \\
\texttt{real\_clips}, \texttt{real\_T} & \Res{tsr.config.cfg.eval.real_clips} / \Res{tsr.config.cfg.eval.real_t} & real test clips per suite, and frames per clip \\
\texttt{eval\_reserve\_h} & \Res{tsr.config.cfg.eval.eval_reserve_h} & hours kept back for the test suites, the evaluation and the export before the training budgets are set \\
\multicolumn{3}{@{}l}{\textit{Pipeline caps (unused, the manifest was not built inline)}} \\
\texttt{list\_cap\_files} & \Res{tsr.config.cfg.pipeline.list_cap_files} & most files listed per Kaggle dataset \\
\texttt{seed}, \texttt{workers} & \Res{tsr.config.cfg.pipeline.seed} / \Res{tsr.config.cfg.pipeline.workers} & seed of the inline pipeline and its worker count \\
\texttt{hf\_caps.game\_stills} & \Res{tsr.config.cfg.pipeline.hf_caps.game_stills.0} / \Res{tsr.config.cfg.pipeline.hf_caps.game_stills.1} & game stills for training and for test \\
\texttt{hf\_caps.tartanair}, \texttt{hf\_caps.tartanair\_seq} & \Res{tsr.config.cfg.pipeline.hf_caps.tartanair} / \Res{tsr.config.cfg.pipeline.hf_caps.tartanair_seq} & frames per TartanAir environment, and frames per sequence \\
\texttt{hf\_caps.gameir} & \Res{tsr.config.cfg.pipeline.hf_caps.gameir.0} / \Res{tsr.config.cfg.pipeline.hf_caps.gameir.1} & GameIR clips for training and for test \\
\texttt{hf\_caps.vimeo1080p} & \Res{tsr.config.cfg.pipeline.hf_caps.vimeo1080p.0} / \Res{tsr.config.cfg.pipeline.hf_caps.vimeo1080p.1} & Vimeo-1080p videos for training and for the pinned validation shard \\
\texttt{calibration\_pairs}, \texttt{val\_fraction} & \Res{tsr.config.cfg.pipeline.calibration_pairs} / \Res{tsr.config.cfg.pipeline.val_fraction} & calibration pairs for the duplicate thresholds, and the validation share \\
\multicolumn{3}{@{}l}{\textit{Training configuration files (typed)}} \\
\texttt{val\_interval\_s} & varies & seconds between validation rounds: 480 for the main runs, and one sixth of the budget for the ablations \\
\texttt{sync\_every} & 10 & steps between the synchronised checks of the clock, of the loss flag and of the learning rate \\
\texttt{progress\_interval\_s} & 15 & seconds between progress lines \\
\texttt{select\_metric} & psnr\_y & validation score that selects the best checkpoint \\
\texttt{rife\_weights} & path & RIFE 4.26 weights loaded into the built-in motion estimator, read from \texttt{/tmp/nss\_rife/rife426.pt} \\
\texttt{teacher\_ckpt} & null & checkpoint of the teacher for distillation, none in TSR \\
\texttt{workers\_per\_rank} & 2 & data-loading worker processes per GPU process, half the CPU cores \\
\texttt{weights\_dir}, \texttt{metrics\_dir} & folders & output folders of the run, under \texttt{/kaggle/working} \\
\end{longtable}



\Cref{tab:app-tsr-runs} lists the TSR training runs as the notebook's launcher builds them. The 2x student is the \tier{low} tier at scale 2, and its cost is \Res{tsr.model_tiers.low.2.gmac_per_frame_1080p} GMAC per 1080p frame. Every run used a batch of 48 sequences per GPU, the largest size of the batch probe that fitted for both the teacher and the student.

\begin{table}[!htbp]
\caption{TSR training runs. Each main run used both GPUs, and each pair ran one run per GPU in parallel. The six ablation runs form three pairs of 30-minute runs, and all use the \tier{low} tier at scale 3.}\label{tab:app-tsr-runs}
\footnotesize
\begin{tabular*}{\textwidth}{@{\extracolsep\fill}llrlrrr}
\toprule
Run & Tier & Scale & GPUs & Batch & Budget (min) & Peak lr \\
\midrule
\texttt{tsr\_teacher\_s3} & \tier{teacher} & \Res{tsr.config.cfg.scale} & both & 48 & \Res{tsr.config.cfg.minutes.teacher} & \Res{tsr.config.cfg.lr_teacher} \\
\texttt{tsr\_low\_s3} & \tier{low} & \Res{tsr.config.cfg.scale} & both & 48 & \Res{tsr.config.cfg.minutes.student} & \Res{tsr.config.cfg.lr} \\
\texttt{tsr\_igpu\_s3} & \tier{igpu} & \Res{tsr.config.cfg.scale} & GPU 0 & 48 & \Res{tsr.config.cfg.minutes.pair} & \Res{tsr.config.cfg.lr} \\
\texttt{tsr\_low\_s2} & \tier{low} & \Res{tsr.config.cfg.scale_alt} & GPU 1 & 48 & \Res{tsr.config.cfg.minutes.pair} & \Res{tsr.config.cfg.lr} \\
\texttt{abl\_*} (six runs) & \tier{low} & \Res{tsr.config.cfg.scale} & one each & 48 & \Res{tsr.config.cfg.minutes.ablation} each & \Res{tsr.config.cfg.lr} \\
\botrule
\end{tabular*}
\end{table}

\Cref{tab:app-tsr-abl} lists the ablations. Each changes only the keys listed, so all six share the reference's other keys and its budget.

\begin{table}[!htbp]
\caption{TSR ablations, from the notebook's \texttt{ABLATIONS} dictionary. Each row lists the keys that differ from the reference run.}\label{tab:app-tsr-abl}
\footnotesize
\begin{tabular*}{\textwidth}{@{\extracolsep\fill}lp{0.30\textwidth}p{0.42\textwidth}}
\toprule
Ablation & Keys that differ & What the run tests \\
\midrule
\texttt{reference} & none & the full recipe at the same short budget \\
\texttt{no\_history} & \texttt{use\_history} false, \texttt{find\_unused} true & no recurrent history, so each frame is upscaled alone. The unused-parameter flag of distributed training is on, because the history branch receives no gradient \\
\texttt{no\_gbuffers} & \texttt{ablate\_gbuffers} true & colour-only input, so the motion comes from the built-in estimator and no engine motion vectors or depth are used \\
\texttt{no\_motion\_net} & \texttt{motion\_net} false, \texttt{ablate\_gbuffers} true & colour-only input and no built-in estimator, so the motion is zero \\
\texttt{no\_video\_training} & \texttt{p\_video\_synth} and \texttt{p\_video\_real} 0 & only render-like samples, with no plain video in training \\
\texttt{no\_long\_unroll} & \texttt{long\_every} 0 & only the six-frame sequences, with no long unroll \\
\botrule
\end{tabular*}
\end{table}

\subsubsection{Frame generation}\label{sec:app-cfg-fg}

The FG training runs use the same optimiser, schedule and training loop as TSR, with their own loss terms, and their distillation term is zero when no teacher checkpoint is given. \Cref{tab:app-cfg-fg} lists the keys of the FG configuration. The keys that differ between the four training runs are in \cref{tab:app-fg-runs}.

\begin{table}[!htbp]
\caption{FG training runs, from \texttt{code/cfg\_fg\_*.json}. The four runs share every other key of \cref{tab:app-cfg-fg}, including the training-only keys. The budgets are the planned budgets, scaled slightly down before training.}\label{tab:app-fg-runs}
\footnotesize
\begin{tabular*}{\textwidth}{@{\extracolsep\fill}lllll}
\toprule
Key & \tier{high} (\texttt{fg\_teacher}) & \tier{low} (\texttt{fg\_low}) & \tier{mid} (\texttt{fg\_mid}) & \tier{igpu} (\texttt{fg\_igpu}) \\
\midrule
\texttt{name} & \texttt{fg\_teacher} & \texttt{fg\_low} & \texttt{fg\_mid} & \texttt{fg\_igpu} \\
\texttt{tier} & teacher & low & mid & igpu \\
\texttt{batch\_per\_gpu} & 48 & 48 & 24 & 48 \\
\texttt{lr} & \Res{fg.config.cfg.lr_teacher} & \Res{fg.config.cfg.lr} & \Res{fg.config.cfg.lr} & \Res{fg.config.cfg.lr} \\
\texttt{teacher\_ckpt} & none & \texttt{fg\_teacher\_best.pt} & \texttt{fg\_teacher\_best.pt} & \texttt{fg\_teacher\_best.pt} \\
\texttt{train\_seconds} & scaled teacher & scaled student & scaled pair & scaled pair \\
\botrule
\end{tabular*}
\end{table}

\Cref{tab:app-fg-tiers} lists the FG tiers. The teacher has the architecture of the \tier{high} tier, which is the RIFE 4.26 network with the engine-hint channels and the residual and UI head added. The smaller tiers inherit the leading channel slices of the RIFE weights, and the \tier{teacher} and \tier{high} rows are the same network.

{\footnotesize\setlength{\LTcapwidth}{\textwidth}
\begin{longtable}{lllrrrrr}
\caption{FG tiers. Flow widths and residual blocks describe the IFNet, refinement is the width of the half-resolution residual and UI head, and flow lines is the largest image height at which the flow is estimated, taller frames being halved. The flow scale is the factor applied to a 1080p frame, and the cost is per 1080p output frame.}\label{tab:app-fg-tiers}\\
\toprule
Tier & Flow widths & Blocks & Refine & Flow lines & Flow scale & Params & GMAC \\
\midrule
\tier{igpu} & 64/48/32/24/16 & 3 & 12 & 576 & \Res{fg.model_tiers.igpu.flow_scale_1080p} & \Res{fg.model_tiers.igpu.params} & \Res{fg.model_tiers.igpu.gmac_per_1080p_frame} \\
\tier{low} & 96/64/48/32/16 & 4 & 16 & 576 & \Res{fg.model_tiers.low.flow_scale_1080p} & \Res{fg.model_tiers.low.params} & \Res{fg.model_tiers.low.gmac_per_1080p_frame} \\
\tier{mid} & 128/96/64/48/24 & 6 & 24 & 1152 & \Res{fg.model_tiers.mid.flow_scale_1080p} & \Res{fg.model_tiers.mid.params} & \Res{fg.model_tiers.mid.gmac_per_1080p_frame} \\
\tier{high} & 192/128/96/64/32 & 8 & 32 & 1152 & \Res{fg.model_tiers.high.flow_scale_1080p} & \Res{fg.model_tiers.high.params} & \Res{fg.model_tiers.high.gmac_per_1080p_frame} \\
\botrule
\end{longtable}
}

\FloatBarrier

\begin{longtable}{@{}p{0.31\textwidth}p{0.12\textwidth}p{0.45\textwidth}@{}}
\caption{FG configuration. Keys are under \texttt{cfg} in \texttt{metrics/config.json} unless shown otherwise.}\label{tab:app-cfg-fg}\\
\toprule
Key & Value & Meaning \\
\midrule
\endfirsthead
\toprule
Key & Value & Meaning \\
\midrule
\endhead
\midrule
\multicolumn{3}{r}{\textit{Continued on the next page}} \\
\endfoot
\botrule
\endlastfoot
\multicolumn{3}{@{}l}{\textit{Budget, model and data}} \\
\texttt{budget\_h} (top level) & \Res{fg.config.budget_h} & ceiling of the run in hours \\
\texttt{seed} & \Res{fg.config.cfg.seed} & seed of the training run \\
\texttt{patch} & \Res{fg.config.cfg.patch} & training crop in pixels. The coarsest RIFE level works at one sixteenth of the crop \\
\texttt{still\_size} & \Res{fg.config.cfg.still_size} & side in pixels of the textures from which synthetic scenes are rendered \\
\texttt{student\_tier}, \texttt{small\_tier}, \texttt{mid\_tier} & low, igpu, mid & tiers of the student, of the integrated-GPU run and of the mid-size run (\cref{tab:app-fg-tiers}). The teacher is the \tier{high} tier. \\
\multicolumn{3}{@{}l}{\textit{Caches (requested counts)}} \\
\texttt{still\_train}, \texttt{still\_val} & \Res{fg.config.cfg.cache.still_train} / \Res{fg.config.cfg.cache.still_val} & still-image crops for training and for validation \\
\texttt{quad\_train}, \texttt{quad\_val} & \Res{fg.config.cfg.cache.quad_train} / \Res{fg.config.cfg.cache.quad_val} & quads, meaning four consecutive frames of a video, for training and for validation \\
\texttt{build\_min} & \Res{fg.config.cfg.cache.build_min} & total minutes allowed for decoding all caches \\
\multicolumn{3}{@{}l}{\textit{Sampling}} \\
\texttt{p\_real} & \Res{fg.config.cfg.p_real} & share of each batch drawn from real quads, the rest being synthetic \\
\texttt{p\_extrap} & \Res{fg.config.cfg.p_extrap} & share of samples that extrapolate a future frame from the two latest frames \\
\texttt{p\_tau} & \Res{fg.config.cfg.p_tau} & share of synthetic interpolation samples whose time is drawn at random between 0.15 and 0.85 \\
\texttt{p\_mfg} & \Res{fg.config.cfg.p_mfg} & share of real interpolation samples taken as 3x multi-frame generation \\
\texttt{p\_drop\_mv}, \texttt{p\_drop\_mv\_real} & \Res{fg.config.cfg.p_drop_mv} / \Res{fg.config.cfg.p_drop_mv_real} & share of synthetic and of real samples whose motion vectors are withheld \\
\texttt{p\_ui}, \texttt{p\_menu} & \Res{fg.config.cfg.p_ui} / \Res{fg.config.cfg.p_menu} & share of synthetic samples with HUD and subtitle overlays, and with full-screen menus \\
\texttt{p\_cut}, \texttt{p\_flip} & \Res{fg.config.cfg.p_cut} / \Res{fg.config.cfg.p_flip} & share of synthetic samples with a scene cut, and with a rapid camera reversal \\
\multicolumn{3}{@{}l}{\textit{Optimisation}} \\
\texttt{lr}, \texttt{lr\_teacher} & \Res{fg.config.cfg.lr} / \Res{fg.config.cfg.lr_teacher} & peak learning rate of the students and of the teacher, followed by warm-up and a cosine decay over the wall-clock budget \\
\texttt{min\_lr}, \texttt{warmup\_frac} & \Res{fg.config.cfg.min_lr} / \Res{fg.config.cfg.warmup_frac} & floor of the cosine decay, and share of the budget spent warming up \\
\texttt{weight\_decay}, \texttt{ema\_decay} & \Res{fg.config.cfg.weight_decay} / \Res{fg.config.cfg.ema_decay} & AdamW weight decay, and decay of the evaluated moving average of the weights \\
\texttt{grad\_clip} & \Res{fg.config.cfg.grad_clip} & gradient norm above which the gradients are scaled down \\
\multicolumn{3}{@{}l}{\textit{Loss weights (per-sample loss terms)}} \\
\texttt{w\_census} & \Res{fg.config.cfg.w_census} & weight of the census loss, a soft ternary distance that tolerates brightness changes \\
\texttt{w\_lpips} & \Res{fg.config.cfg.w_lpips} & weight of the LPIPS perceptual loss, set to zero when the package is missing \\
\texttt{w\_flow} & \Res{fg.config.cfg.w_flow} & weight of the flow supervision, per pixel of flow error, with RAFT labels on video at a quarter of this weight \\
\texttt{w\_ui} & \Res{fg.config.cfg.w_ui} & weight of the cross-entropy on the UI mask \\
\texttt{w\_distill} & \Res{fg.config.cfg.w_distill} & weight of distillation of the flows and output from the teacher, zero without a teacher checkpoint \\
\multicolumn{3}{@{}l}{\textit{Validation and budgets}} \\
\texttt{val\_bs}, \texttt{val\_batches} & \Res{fg.config.cfg.val_bs} / \Res{fg.config.cfg.val_batches} & samples per validation batch and batches per validation round \\
\texttt{val\_interval\_min} & \Res{fg.config.cfg.val_interval_min} & longest time between validation rounds, in minutes \\
\texttt{minutes.teacher}, \texttt{minutes.student} & \Res{fg.config.cfg.minutes.teacher} / \Res{fg.config.cfg.minutes.student} & planned wall-clock budget of the teacher and of the student, in minutes \\
\texttt{minutes.pair}, \texttt{minutes.ablation} & \Res{fg.config.cfg.minutes.pair} / \Res{fg.config.cfg.minutes.ablation} & planned budget of each parallel pair, and of each ablation pair, in minutes \\
\multicolumn{3}{@{}l}{\textit{Evaluation}} \\
\texttt{synthetic\_per\_suite} & \Res{fg.config.cfg.eval.synthetic_per_suite} & synthetic test crops per synthetic suite \\
\texttt{real\_clips}, \texttt{native\_clips} & \Res{fg.config.cfg.eval.real_clips} / \Res{fg.config.cfg.eval.native_clips} & real clips per real suite, and clips of the native 1080p suite \\
\texttt{eval\_reserve\_h} & \Res{fg.config.cfg.eval.eval_reserve_h} & hours kept back for the test suites, the evaluation and the export before the training budgets are set \\
\multicolumn{3}{@{}l}{\textit{Pipeline caps (unused, the manifest was not built inline)}} \\
\texttt{list\_cap\_files} & \Res{fg.config.cfg.pipeline.list_cap_files} & most files listed per Kaggle dataset \\
\texttt{seed}, \texttt{workers} & \Res{fg.config.cfg.pipeline.seed} / \Res{fg.config.cfg.pipeline.workers} & seed of the inline pipeline and its worker count \\
\texttt{hf\_caps.game\_stills} & \Res{fg.config.cfg.pipeline.hf_caps.game_stills.0} / \Res{fg.config.cfg.pipeline.hf_caps.game_stills.1} & game stills for training and for test \\
\texttt{hf\_caps.tartanair}, \texttt{hf\_caps.tartanair\_seq} & \Res{fg.config.cfg.pipeline.hf_caps.tartanair} / \Res{fg.config.cfg.pipeline.hf_caps.tartanair_seq} & frames per TartanAir environment, and frames per sequence \\
\texttt{hf\_caps.gameir} & \Res{fg.config.cfg.pipeline.hf_caps.gameir.0} / \Res{fg.config.cfg.pipeline.hf_caps.gameir.1} & GameIR clips, both zero, so GameIR was not used \\
\texttt{hf\_caps.vimeo1080p} & \Res{fg.config.cfg.pipeline.hf_caps.vimeo1080p.0} / \Res{fg.config.cfg.pipeline.hf_caps.vimeo1080p.1} & Vimeo-1080p videos for training and for the pinned validation shard \\
\texttt{calibration\_pairs}, \texttt{val\_fraction} & \Res{fg.config.cfg.pipeline.calibration_pairs} / \Res{fg.config.cfg.pipeline.val_fraction} & calibration pairs for the duplicate thresholds, and the validation share \\
\multicolumn{3}{@{}l}{\textit{Training configuration files (typed)}} \\
\texttt{val\_interval\_s} & varies & seconds between validation rounds: 480 for the four main runs, and one sixth of the budget for the ablations \\
\texttt{sync\_every} & 10 & steps between the synchronised checks of the clock, of the loss flag and of the learning rate \\
\texttt{progress\_interval\_s} & 15 & seconds between progress lines \\
\texttt{select\_metric} & psnr & validation score that selects the best checkpoint \\
\texttt{rife\_weights} & path & RIFE 4.26 weights loaded into the IFNet, unless the run is the no-inheritance ablation, read from \texttt{/tmp/nss\_rife/rife426.pt} \\
\texttt{ui\_training} & true & whether the UI head is trained, false only in the no-UI-training ablation \\
\texttt{workers\_per\_rank} & 2 & data-loading worker processes per GPU process, half the CPU cores \\
\texttt{weights\_dir}, \texttt{metrics\_dir} & folders & output folders of the run, under \texttt{/kaggle/working} \\
\end{longtable}


\FloatBarrier


\Cref{tab:app-fg-abl} lists the FG ablations. Each changes only the keys listed. The ablations run for the ablation budget, and their validation interval is one sixth of that budget.

\begin{table}[!htbp]
\caption{FG ablations, from \texttt{code/cfg\_abl\_*.json}. Each row lists the keys that differ from the reference run, which is the \tier{low} tier with the full recipe.}\label{tab:app-fg-abl}
\footnotesize
\begin{tabular*}{\textwidth}{@{\extracolsep\fill}lp{0.30\textwidth}p{0.42\textwidth}}
\toprule
Ablation & Keys that differ & What the run tests \\
\midrule
\texttt{reference} & none & the full recipe at the same short budget \\
\texttt{no\_ui\_training} & \texttt{ui\_training} false, \texttt{w\_ui} 0 & no HUD, subtitle or menu overlays in training, and no UI loss \\
\texttt{no\_motion\_prior} & \texttt{ablate\_prior} true & no engine motion vectors and no camera prior \\
\texttt{extrap\_only} & \texttt{mode} extrap & every sample extrapolates, with no interpolation \\
\texttt{no\_distillation} & \texttt{teacher\_ckpt} null & no teacher, so the distillation term is zero \\
\texttt{no\_inheritance} & \texttt{rife\_weights} null & random initialisation instead of the RIFE 4.26 channel slices \\
\botrule
\end{tabular*}
\end{table}

\subsection{Reproducibility}\label{sec:app-repro}

The steps below are in the order that a new run needs them. The notebooks are built from the repository's sources and run on Kaggle. Their saved outputs are read by two scripts, which write the values that the text quotes and both documents.

\subsubsection{Building the notebooks}\label{sec:app-build}

Each \texttt{notebooks/<name>.py} file is a percent-format script, and its cells become the cells of the notebook. A cell whose only line is \texttt{\# WRITEFILE} followed by a file name of \texttt{src/neuralss} becomes a cell that writes that file into \texttt{/kaggle/working/code}. Every notebook therefore carries its own copy of the library, and the library has a single source. \texttt{tools/build\_notebooks.py} writes the three \texttt{.ipynb} files with stable cell identifiers, and it rewrites a notebook only when its sources changed, which removes that notebook's outputs. Run it from the repository root with the repository's Python in isolated mode (\texttt{python -I tools/build\_notebooks.py}). Its \texttt{--check} option exits with code 1 when a notebook does not match its sources, and ignores outputs. The build needs \texttt{nbformat}.

\subsubsection{Rerunning on Kaggle}\label{sec:app-kaggle}

Each notebook runs on Kaggle with the accelerator set to GPU T4 x2 and Internet switched on. The notebooks download the RIFE 4.26 weights, checked against SHA-256 hashes, and the pre-trained RAFT weights, and they read Hugging Face data for any source that is not attached. The optional secret \texttt{HF\_TOKEN} raises the Hugging Face download limits. The hardware check stops the notebook unless it sees two GPUs.

\begin{enumerate}
\item Create the data pipeline notebook from \texttt{neural-data-pipeline/\allowbreak neural-data-pipeline.ipynb}, attach the seven Kaggle datasets listed in its header (DIV2K, Flickr2K, REDS, Vimeo-90K septuplets and triplets, MPI Sintel and Vid4), and choose \emph{Save \& Run All}.
\item Create the TSR and FG notebooks from their \texttt{.ipynb} files, attach the same datasets, attach the output of the data pipeline run with \emph{Add Input}, and choose \emph{Save \& Run All}. Both read the manifest from that output.
\item A full run takes about the budget of its notebook. To check a change first, set \texttt{QUICK\_RUN = True} in the configuration cell, which the notebooks estimate at about 40 minutes.
\item If Kaggle reports \texttt{ERRORED\_MOUNTING\_DATASET}, the session never starts and shows no log. Remove or re-add the dataset that failed, and save again.
\end{enumerate}

\subsubsection{Where the outputs go}\label{sec:app-outputs}

Each notebook writes to \texttt{/kaggle/working}. The training notebooks zip everything in that folder except the training caches, which they delete first, and the data pipeline zips its metrics, plots, logs, code and manifest. The \texttt{outputs} folder of each notebook in the repository has the same layout: \texttt{code} holds the library and the run configurations as written on Kaggle, \texttt{logs} holds the notebook log and one log per training run, \texttt{metrics} holds every CSV and JSON result, \texttt{plots} holds every figure and, for the training notebooks, \texttt{predictions} and \texttt{weights} hold the sample frames and the checkpoints. To update a notebook, extract its output into \texttt{<notebook>/outputs}, and save the Kaggle copy of the notebook, with its outputs, as \texttt{outputs/<name>.executed.ipynb}. The generator reads each run's dates from that copy if it exists and otherwise from the notebook file itself. The data pipeline has no such copy in \texttt{outputs}, because its executed notebook is the notebook file, so that file must keep its outputs.

\subsubsection{Regenerating the numbers and the documents}\label{sec:app-generate}

\texttt{tools/build\_paper\_results.py} reads the metric files of the three \texttt{outputs/metrics} folders, the run logs, the executed notebooks and the Middlebury comparison file. It writes \texttt{paper/results-values.tex}, which holds one \texttt{\textbackslash SetResult} line per key. A metric file that has no rule stops the script, so that no number is left out without notice. The script scans the \texttt{paper} sources for \texttt{\textbackslash Res} uses and prints a change report of the added, removed and changed keys, with the sentences that use each removed or changed key. A use of a key that does not exist is reported with its file and line, and the script exits with code 1 after writing the file. The script also copies the five \texttt{fig-tsr-*.png} figures from their notebook outputs when their bytes differ. The six \texttt{fig-fg-ui-*.png} figures are not regenerated, because each is a resampled copy of a notebook image and the resampling cannot be reproduced exactly.

\texttt{tools/build\_paper.py} runs the generator and then compiles the six documents with Tectonic in \texttt{paper}. The article and its supplement are built in the order \texttt{main}, \texttt{supplementary}, \texttt{main}, so that each one's cross-references resolve against the other. The anonymised pair follows, then the title page and the copy without references. It prints the page count of each PDF. The option \texttt{--no-generate} skips the generator, and \texttt{--supplement-only} compiles the supplement alone. Tectonic is \texttt{paper/tectonic.exe} when that file exists, and otherwise \texttt{tectonic} on the \texttt{PATH}.

\subsubsection{The Middlebury comparison}\label{sec:app-mb}

The Middlebury comparison of the FG tiers is not produced by the notebooks. \texttt{tools/compare\_fg\_tiers\_middlebury.py} runs on a CPU and needs three inputs, each of which can be given on the command line. The first is the \texttt{other-color-allframes} and \texttt{other-gt-interp} archives of the Middlebury optical-flow benchmark \cite{baker2011middlebury}, extracted into one folder. The second is the RIFE 4.26 \texttt{flownet.pkl} from the \texttt{RIFEv4.26\_0921.zip} archive of \texttt{huggingface.co/hzwer/RIFE}, whose SHA-256 hash the notebooks check when they download it. The third is the four \texttt{fg\_*\_best.pt} checkpoints of the FG run. The script defaults expect the downloads under \texttt{.agent-local/downloads}, a folder that git does not track, so pass the paths explicitly. It evaluates frame averaging, RIFE 4.26 at full, half and quarter flow resolution, and the four FG tiers, on the twelve sequences. It writes \texttt{neural-frame-generation/local-evaluation/middlebury\_psnr.csv}, with one row per sequence and method and the mean of each method with its cost in GMAC per 1080p frame. The generator reads that file as the \texttt{mb} keys, so rerun the generator after the script.

\subsubsection{Weights and exported models}\label{sec:app-weights}

Each training run saves three PyTorch checkpoints in its \texttt{weights} folder. The file \texttt{<run>\_best.pt} holds the exponential-moving-average weights of the best validation round. The file \texttt{<run>\_last.pt} holds the weights of the last round, and \texttt{<run>\_fused.pt} holds the network with its RepConv branches folded into plain convolutions. \Cref{tab:app-weights} lists the models that the two training runs exported. The checkpoints are research weights, trained on research-use data, and the repository's \texttt{RESEARCH.md} records the data licences. The notebooks load RIFE 4.26 \cite{huang2022rife} into the motion estimators and the FG network, and they log how many tensors were copied exactly, copied by slicing or left missing. No run left a tensor missing.

\begin{table}[!htbp]
\caption{Weights and exported models of the two training runs. Each name stands for the three checkpoints of its run (\texttt{\_best}, \texttt{\_last} and \texttt{\_fused}) unless it names an ONNX file. File sizes are in the repository's README.}\label{tab:app-weights}
\footnotesize
\begin{tabular*}{\textwidth}{@{\extracolsep\fill}p{0.42\textwidth}lp{0.40\textwidth}}
\toprule
Files & Notebook & Model \\
\midrule
\texttt{tsr\_low\_s3\_*.pt} & TSR & main 3x student, \tier{low} tier \\
\texttt{tsr\_igpu\_s3\_*.pt} & TSR & integrated-GPU 3x student \\
\texttt{tsr\_low\_s2\_*.pt} & TSR & 2x student, \tier{low} tier \\
\texttt{tsr\_teacher\_s3\_*.pt} & TSR & 3x teacher, the quality ceiling of the architecture \\
\texttt{abl\_*\_*.pt} & TSR & the six ablation runs \\
\texttt{tsr\_student\_s3\_fused.onnx}, \texttt{tsr\_igpu\_s3\_fused.onnx} & TSR & one recurrent step of the 3x student and of the integrated-GPU student, with static shapes \\
\texttt{fg\_teacher\_*.pt} & FG & \tier{high} tier, RIFE 4.26 fine-tuned, the strongest FG model \\
\texttt{fg\_mid\_*.pt}, \texttt{fg\_low\_*.pt}, \texttt{fg\_igpu\_*.pt} & FG & \tier{mid}, \tier{low} (the student) and \tier{igpu} tiers \\
\texttt{abl\_*\_*.pt} & FG & the six ablation runs \\
\texttt{fg\_student\_fused.onnx}, \texttt{fg\_igpu\_fused.onnx} & FG & \tier{low} (student) and \tier{igpu} tiers at 128 by 128 pixels, with static shapes \\
\botrule
\end{tabular*}
\end{table}

Each ONNX file was exported with the TorchScript exporter at opset 17 and checked against PyTorch with ONNX Runtime on the CPU. The maximum difference of each check is in \texttt{onnx\_export.json}. The ONNX files have static shapes, so deployment at another resolution needs a new export with \texttt{nss\_models.export\_onnx}.

\subsection{Revising the supplement after a new run}\label{sec:app-revise}

A new run changes the numbers, and some of the sentences that quote them. The following procedure keeps the text consistent with the saved outputs.

\begin{enumerate}
\item Edit the notebook sources in \texttt{notebooks} or the library in \texttt{src/neuralss}, never the \texttt{.ipynb} files. Rebuild the notebooks, check them with \texttt{--check}, and run the lint and format checks that the repository's README lists.
\item Store the executed copy of each notebook that is being replaced, with its outputs, as described in \cref{sec:app-outputs}, before the rebuild removes its outputs.
\item Run the notebook on Kaggle as in \cref{sec:app-kaggle}, and replace its \texttt{outputs} folder with the new output.
\item Rerun \texttt{tools/build\_paper\_results.py} and read its change report. Each removed or changed key lists the sentences that use it, and each of them needs checking against the new value and its meaning.
\item Revisit each listed sentence and its status label, as defined in the first section of this supplement. A finding whose value moves by more than the run-to-run variation described there needs its label checked, and its sentence must describe the new value, not only quote its number.
\item Update the run descriptions in the notebook headers and in the repository's README, which quote earlier runs. Update the typed values of \cref{tab:app-tsr-runs,tab:app-fg-runs} and of the training-only rows from the new \texttt{code/cfg\_*.json} files and notebook sources, because the generator does not read them.
\item If the checkpoints changed, rerun \texttt{tools/compare\_fg\_tiers\_middlebury.py}, then the generator, since the Middlebury values are not produced on Kaggle.
\item Rebuild the documents with \texttt{tools/build\_paper.py}, read the page counts it prints, and check the article against its twelve-page limit before the supplement is finalised. Then view every rendered page.
\end{enumerate}
